\documentclass[11pt]{article}
\usepackage[T1]{fontenc}
\usepackage{lmodern}
\usepackage[letterpaper,hmargin=1.02in,vmargin=0.95in]{geometry}
\usepackage{mathtools,amssymb,amsthm,microtype,xcolor,needspace}
\definecolor{inkblue}{RGB}{25,66,85}
\usepackage[colorlinks=true,linkcolor=inkblue,citecolor=inkblue,urlcolor=inkblue]{hyperref}
\newtheorem{theorem}{Theorem}
\newtheorem{lemma}{Lemma}
\newtheorem{proposition}{Proposition}
\newtheorem{corollary}{Corollary}
\newcommand{\e}{\mathrm e}
\newcommand{\Q}{\mathbb Q}
\newcommand{\N}{\mathbb N}
\newcommand{\M}{\mathcal M}
\DeclareMathOperator{\den}{den}
\DeclareMathOperator{\lcm}{lcm}
\setlength{\parskip}{3pt}
\setlength{\emergencystretch}{2em}
\hypersetup{pdftitle={A New Family of Normal Numbers},
pdfsubject={The Xi family, its Stoneham specialization, and normality of finite rational linear combinations}}

\begin{document}
\begin{center}
{\LARGE\bfseries\color{inkblue} A New Family of Normal Numbers\par}
\vspace{9pt}
{\small 10 September 2026\par}
\end{center}

\begin{abstract}
We study the numbers \(\xi_{b,c;a}\), obtained by restricting the
logarithmic series to products of powers of \(a\) and \(c\). For
integers \(a,b,c\ge2\) with \(\gcd(c,ab)=1\), we prove normality in
base \(b\). More strongly, fixing \(a,b\) and varying the eligible
parameter \(c\), every finite rational combination of distinct members
with coefficients not all zero is normal, even after a rational translation. The
construction extends Stoneham numbers and includes values of Dibag's
localized logarithms. The proof follows the denominators of rational
approximations: one prime detects a contribution that the others cannot
cancel. Both main theorems have complete Lean formalizations.
\end{abstract}

\section{Meet the Xi family}

A real number is \emph{normal in base \(b\)} if every word of \(\ell\)
base-\(b\) digits occurs with limiting frequency \(b^{-\ell}\), counting
overlapping occurrences. Thus a binary normal number contains each digit
half the time, each two-digit word a quarter of the time, and so on.
For numbers outside \([0,1)\), we use the expansion of the fractional part.

The numbers in this paper have a simple analytic definition. For positive
integers \(a,c\), write
\[
 \M(a,c)=\{a^i c^k:i,k\ge0\},
\]
with each integer included only once. For an integer \(b\ge2\), define
\begin{equation}\label{eq:xi}
 \boxed{\displaystyle
 \xi_{b,c;a}=\sum_{m\in\M(a,c)}\frac1{m b^m}.}
\end{equation}
We call these numbers the \emph{Xi family}. The first index is the
expansion base; the index after the semicolon is an additional
multiplicative generator.

The series converges, since it is a subseries of
\[
 \sum_{m\ge1}\frac1{m b^m}=-\log(1-1/b).
\]
When \(a,c\ge2\) are coprime, every product \(a^ic^k\) has a unique
representation, so \eqref{eq:xi} becomes the double series
\begin{equation}\label{eq:double}
 \xi_{b,c;a}=\sum_{i,k\ge0}\frac1{a^i c^k b^{a^i c^k}}.
\end{equation}
Indeed, an equality between two such products would equate a positive
power of \(a\) with a positive power of \(c\), unless both pairs of
exponents agree. Coprimality excludes that possibility.

For example, \(\xi_{2,3;2}\) uses the entries of the grid
\[
\begin{array}{c|rrrrr}
 & 1&3&9&27&\cdots\\ \hline
1&1&3&9&27&\cdots\\
2&2&6&18&54&\cdots\\
4&4&12&36&108&\cdots\\
8&8&24&72&216&\cdots\\
\vdots&\vdots&\vdots&\vdots&\vdots&
\end{array}
\qquad\text{each entry \(m\) contributes \(1/(m2^m)\).}
\]
Its value begins \(0.68562264189340237754\ldots\). The rule selects
terms of a familiar convergent series; it makes no reference to digit
frequencies.

\begin{theorem}[Normality of the Xi family]\label{thm:normal}
Let \(a,b,c\ge2\) be integers with \(\gcd(c,ab)=1\). Then
\(\xi_{b,c;a}\) is normal in base \(b\).
\end{theorem}

The condition allows \(a\) to share prime factors with the base. It
covers \(\xi_{2,3;2}\), for instance, and also the composite-generator
example \(\xi_{2,35;6}\).

The arithmetic property is stronger. Addition can destroy normality:
if \(x\in(0,1)\) is normal, so is \(1-x\), but their sum is an
integer. Within the following fixed-parameter Xi family, rational
combinations behave differently.

\begin{theorem}[Rational combinations]\label{thm:arithmetic}
Fix integers \(a,b\ge2\). Let \(c_1,\ldots,c_t\ge2\) be distinct
integers with \(\gcd(c_j,ab)=1\), and let
\(q_0,\lambda_1,\ldots,\lambda_t\in\Q\). If at least one
\(\lambda_j\) is nonzero, then
\begin{equation}\label{eq:main-combination}
 q_0+\lambda_1\xi_{b,c_1;a}+\cdots+\lambda_t\xi_{b,c_t;a}
\end{equation}
is normal in base \(b\).
\end{theorem}

The parameters may share primes or be powers of one another. In
particular, the binary-normal examples include
\[
 \xi_{2,3;2}-\xi_{2,9;2}
 \quad\text{and}\quad
 \xi_{2,3;2}-2\xi_{2,9;2}+\tfrac17\xi_{2,15;2}.
\]
The first example removes every term whose exponent of three is even.
An infinite portion of the series cancels, and the remainder is still
normal. Accounting for such cancellations is the main work of the proof.

\subsection*{Stoneham numbers and localized logarithms}

The usual Stoneham constant is
\[
 \alpha_{b,c}=\sum_{k\ge1}\frac1{c^k b^{c^k}}.
\]
Stoneham proved foundational normality results for these series;
Bailey and Crandall established base-\(b\) normality for every coprime
pair \(b,c\ge2\) \cite{Stoneham,BaileyCrandall}. Setting the additional
generator to one in our distinct-index definition gives
\begin{equation}\label{eq:stoneham}
 \xi_{b,c;1}=\frac1b+\alpha_{b,c}.
\end{equation}
Thus the construction passes from one geometric sequence of exponents
to a grid of products of powers. Formula~\eqref{eq:stoneham} uses
\eqref{eq:xi}: substituting \(a=1\) in the double series would repeat
every term infinitely. The boundary \(a=1\) is the classical Stoneham
case; our two theorems concern \(a\ge2\).

There is also an older analytic origin. Dibag's localized logarithm
retains the terms of \(\log(1+z)\) whose exponents have prime factors
in a chosen set \cite[Definitions 1.1--1.2]{Dibag}. For distinct prime
generators \(a,c\), our function
\[
 F_{a,c}(z)=\sum_{m\in\M(a,c)}\frac{z^m}{m}
\]
is precisely \(-L_{\{a,c\}}(-z)\) in Dibag's notation, and
\(\xi_{b,c;a}=F_{a,c}(1/b)\). Composite generators allow further
restrictions on the exponents. The results here establish normality
and rational-combination stability for this family.

Our argument belongs to the arithmetic approach to normal numbers
developed by Stoneham, Korobov, Kano, Bailey and Crandall, and others.
Earlier rings of normal numbers provide precedents for arithmetic
closure \cite{Kano,KanoShiokawa}. We use a rational-orbit estimate of
Vandehey \cite{Vandehey}; the accompanying Lean proof independently
establishes the estimates it needs by repeated differencing.

\section{Read the digits through rational approximations}

Multiplication by \(b^n\) moves the digit at position \(n+1\) to the
front. A standard form of Weyl's criterion says that \(x\) is normal
in base \(b\) exactly when
\begin{equation}\label{eq:weyl}
 \sum_{n<N}\e(hb^n x)=o(N)
 \quad\text{for every integer \(h\ne0\)},\qquad
 \e(u)=\exp(2\pi i u).
\end{equation}
Indeed, uniform distribution of \(\{b^nx\}\) gives each interval
\([j/b^\ell,(j+1)/b^\ell)\) its length as frequency. These are exactly
the intervals encoding digit words \cite{Weyl}.

Here is the analytic input. For fixed \(b\), a finite set \(P\) of
primes not dividing \(b\), and any \(\varepsilon>0\), there are
constants \(C,\delta>0\) such that
\begin{equation}\label{eq:orbit}
 \left|\sum_{r=0}^{L-1}\e(Ab^r/q)\right|
 \le C L q^{-\delta}(1+\log q),
 \qquad L\ge q^\varepsilon,
\end{equation}
whenever \(A/q\) is reduced and every prime factor of \(q\) lies in
\(P\). This is Vandehey's Corollary~1.2 \cite{Vandehey}, with the
starting index shifted. The estimate says that a sufficiently long
piece of a rational radix orbit becomes equidistributed as its
denominator grows. All logarithms below are natural.

We will use the following elementary consequence. It is stated for
weighted series because subtracting Xi numbers naturally introduces
weights.

\begin{proposition}[A denominator test]\label{prop:transfer}
Fix \(a,b\ge2\) and integers \(d_1,\ldots,d_g\ge2\) coprime to
\(ab\). Let each \(w_j(k)\) be a bounded rational sequence, with a
common denominator for all its values. Put
\begin{align}
 x&=q_0+\sum_{j=1}^g\sum_{i,k\ge0}
       \frac{w_j(k)}{a^i d_j^k b^{a^i d_j^k}},\label{eq:weighted-x}\\
 T_n&=q_0+\sum_{j=1}^g\sum_{a^i d_j^k\le n}
       \frac{w_j(k)}{a^i d_j^k b^{a^i d_j^k}},
       \qquad q_0\in\Q.\label{eq:weighted-truncation}
\end{align}
Suppose that, for some prime \(p\nmid b\) and some \(\eta>0\),
the reduced denominator of \(T_n\) contains a factor at least
\(n^\eta\) that is a power of \(p\), except at a set of positions
of density zero. Then \(x\) is normal in base \(b\).
\end{proposition}

A set has \emph{density zero} if its intersection with
\(\{1,\ldots,N\}\) has \(o(N)\) elements. Thus occasional denominator
cancellation is allowed; it must occupy a vanishing fraction of digit
positions.

\begin{proof}
Set \(R_n=b^nT_n\). Boundedness of the weights and uniqueness of the
products within each group give
\begin{equation}\label{eq:tail}
 |b^n x-R_n|
 \ll \sum_{m>n}\frac{b^{n-m}}m
 \le\frac1{(b-1)(n+1)}.
\end{equation}
Here \(F\ll G\) means \(|F|\le C_0G\) for a constant depending only
on the fixed data. All implied constants in this proof have this
meaning, with additional dependence on a fixed Fourier frequency allowed.

Let \(D\) clear the coefficient denominators, including that of \(q_0\).
Every denominator of \(R_n\) divides
\begin{equation}\label{eq:den-upper}
 D a^{\lfloor\log_a n\rfloor}
       \prod_{j=1}^g d_j^{\lfloor\log_{d_j}n\rfloor}
 \le D n^{g+1}.
\end{equation}
There are only \(O((\log N)^2)\) insertion positions
\(a^i d_j^k\le2N\). Between successive insertions,
\(R_{n+r}=b^rR_n\) exactly.

Some factors of \(a\) or \(D\) may divide \(b\). On a dyadic
interval \([N,2N)\), remove each insertion position and the following
\(C_1\log(2N)\) positions, including any such intervals crossing \(N\).
For sufficiently large fixed \(C_1\), every denominator prime dividing
\(b\) has then been cleared by multiplication by \(b\). To see this
term by term, its possible denominator exponent at a prime dividing
\(b\) is \(O(\log N)\), whereas each new radix shift supplies a
fixed positive power of that prime. The fixed rational term clears too.
The removed positions number \(O((\log N)^3)\).

The remaining positions form \(O((\log N)^2)\) intervals on each of
which the reduced denominator is constant, coprime to \(b\), supported
on a fixed finite set of primes, and at most \(DN^{g+1}\) up to a
fixed factor. Our hypothesis removes another \(o(N)\) positions:
discard the intervals whose \(p\)-part is smaller than \(N^\eta\).
Each position in such an interval is exceptional, because \(p\nmid b\)
and the \(p\)-part is constant there.

Fix \(h\ne0\). Multiplication by \(h\) reduces a denominator by at
most \(|h|\). Thus on each retained interval, the denominator \(q\)
of \(hR_n\) satisfies
\[
 q\gg_h N^\eta,\qquad q\ll N^{g+1}.
\]
Discard intervals shorter than \(N^{1/2}\); their total length is
\(O(N^{1/2}(\log N)^2)=o(N)\). Choose
\(0<\varepsilon<1/(2(g+1))\). Every remaining interval has length
at least \(q^\varepsilon\) for large \(N\), so \eqref{eq:orbit}
bounds its Fourier sum by its length times \(O_h(N^{-\eta\delta}\log N)\).
Adding the intervals gives
\[
 \sum_{N\le n<M}\e(hR_n)=o(N),\qquad N<M\le2N,
\]
uniformly in \(M\); a final truncated interval is handled by the same
short-interval rule. By \eqref{eq:tail} and
\(|\e(u)-\e(v)|\le2\pi|u-v|\), replacing \(R_n\) by \(b^nx\)
costs only \(O_h(1)\) on this scale. Summing dyadic intervals proves
\eqref{eq:weyl} through every prefix length.
\end{proof}

The analytic task is now complete. To prove normality, it remains to
find one prime whose denominator powers survive most of the time.

\section{Why a denominator survives}

For a prime \(p\), let \(v_p(y)\) be its exponent in a nonzero
rational number \(y\), and put \(v_p(0)=+\infty\). A negative
valuation records a denominator: for example, \(v_3(2/27)=-3\).
We will repeatedly use the elementary rule
\begin{equation}\label{eq:valuation-rule}
 v_p(u+v)=\min(v_p(u),v_p(v))
 \quad\text{if }v_p(u)\ne v_p(v).
\end{equation}
More generally, a finite sum with a unique term of smallest valuation
has that same valuation.

Fix \(a,b,d\ge2\) with \(\gcd(d,ab)=1\), a nonzero periodic
rational sequence \(w(k)\), and a prime \(p\mid d\). Write
\begin{equation}\label{eq:block}
 S_n=\sum_{a^i d^k\le n}\frac{w(k)}{a^i d^k b^{a^i d^k}},
 \qquad D_p(n)=\max\{0,-v_p(S_n)\},
\end{equation}
where \(D_p(n)=0\) when \(S_n=0\). Thus \(D_p(n)\) is the exponent
of \(p\) in the reduced denominator of the whole truncation.

\begin{lemma}[The denominator slope]\label{lem:slope}
With the preceding notation,
\begin{equation}\label{eq:slope}
 \frac{D_p(n)}{\log n}\longrightarrow
 \sigma_p(d):=\frac{v_p(d)}{\log d}
 \quad\text{in density}.
\end{equation}
That is, for every positive tolerance, the positions at which the two
sides differ by more than that tolerance have density zero.
\end{lemma}

The slope has a simple meaning. Among the terms already inserted by
position \(n\), the largest possible exponent of \(d\) is about
\(\log_d n\). It contributes about
\(v_p(d)\log_d n\) powers of \(p\) to the denominator. The lemma
says that cancellation changes this leading size only rarely.

We first record the spacing fact that makes the argument work.

\begin{lemma}[Small relative gaps]\label{lem:gaps}
If \(u,v\ge2\) have no equal positive powers, then for every
\(\varepsilon>0\), every sufficiently large interval
\((x,(1+\varepsilon)x]\) contains a number \(u^i v^k\) with
\(i,k\ge0\).
\end{lemma}

\begin{proof}
The ratio \(\log u/\log v\) is irrational. Its integer multiples
are dense modulo one. Choose finitely many nonnegative multiples whose
fractional parts have gaps smaller than
\(\log(1+\varepsilon)/\log v\). For any sufficiently large
\(\log x\), one of the corresponding values \(i\log u\), followed
by a nonnegative multiple of \(\log v\), lies in
\((\log x,\log x+\log(1+\varepsilon)]\). Exponentiate.
\end{proof}

\begin{proof}[Proof of Lemma~\ref{lem:slope}]
Put \(v=v_p(d)\) and \(\sigma=v/\log d\). Every nonzero summand
at position \(a^id^k\) has valuation
\begin{equation}\label{eq:term-valuation}
 v_p(w(k))-kv,
\end{equation}
since \(p\nmid ab\). There are only finitely many nonzero weights,
so \(k\le\log_d n\) immediately gives
\[
 D_p(n)\le\sigma\log n+O(1).
\]
The lower bound requires three observations.

\paragraph{Equal valuations cannot occur too close together.}
If two expressions in \eqref{eq:term-valuation} agree, the difference
of their \(k\)-coordinates belongs to a fixed finite set. The ratio
of the two insertion positions is therefore \(a^r d^s\), where
\(r\in\mathbb Z\) and \(s\) lies in that finite set. For each
fixed \(s\), the logarithms \(r\log a+s\log d\) form a discrete
set. After excluding zero, their distances from zero have a positive
lower bound, uniformly over these finitely many \(s\).

Choose \(\Lambda>1\) smaller than the resulting separation ratio.
Distinct nonzero terms inserted in any interval \([X,\Lambda X]\)
then have distinct valuations. Any nonempty sum of such terms has the
valuation of its unique smallest-valuation term.

\paragraph{Deep denominator terms occur frequently on a relative scale.}
Fix \(0<\theta<\sigma\). Let \(L\ge1\) be a period of \(w\),
and choose \(k_0\) with \(w(k_0)\ne0\). At scale \(N\), choose
the least nonnegative \(k_N\equiv k_0\pmod L\) for which
\[
 k_Nv-v_p(w(k_0))>\theta\log N.
\]
Then
\begin{equation}\label{eq:threshold}
 d^{k_N}\asymp N^{\theta/\sigma}.
\end{equation}
Every position in the set
\begin{equation}\label{eq:deep-set}
 d^{k_N}\M(a,d^L)
\end{equation}
carries a nonzero term of valuation smaller than
\(-\theta\log N\). Here \(F\asymp G\) means that their ratio is
bounded above and below by positive fixed constants.

The generators \(a,d^L\) are coprime and hence have no equal positive
powers. By Lemma~\ref{lem:gaps}, the set in \eqref{eq:deep-set}
meets every interval \((x,(1+\varepsilon)x]\), uniformly for
\(N\le x\le2N\), once \(N\) is sufficiently large. Indeed,
after division by \(d^{k_N}\), the starting points are at least a
constant times \(N^{1-\theta/\sigma}\), which tends to infinity.

\paragraph{Extensive cancellation is confined to short intervals.}
Call a position \(n\in[N,2N]\) bad if
\(D_p(n)\le\theta\log N\), equivalently
\(v_p(S_n)\ge-\theta\log N\). Take two bad positions \(x<y\)
within a block of multiplicative ratio at most \(\Lambda\).
Their difference \(S_y-S_x\) cannot contain an insertion from
\eqref{eq:deep-set}: the inserted terms have distinct valuations,
so such an insertion would force
\(v_p(S_y-S_x)<-\theta\log N\), contradicting the bound at both
endpoints.

But a deep insertion occurs before \((1+\varepsilon)x\).
Consequently two bad positions in this block are separated by less
than \(\varepsilon x\le2\varepsilon N\). All bad positions in the
block therefore lie in an interval of length at most
\(2\varepsilon N\). Cover \([N,2N]\) by a fixed finite number of
such blocks. Since \(\varepsilon\) is arbitrary,
\begin{equation}\label{eq:bad-count}
 \#\{n\in[N,2N]:D_p(n)\le\theta\log N\}=o(N).
\end{equation}
This holds for every \(\theta<\sigma\). To replace \(\log N\)
by \(\log n\), first increase \(\theta\) slightly while keeping
it below \(\sigma\); on this scale the two logarithms differ by at
most \(\log2\). Together with the upper bound, and then a dyadic
summation, this proves \eqref{eq:slope}.
\end{proof}

\section{Normality of each Xi number}

We can now prove Theorem~\ref{thm:normal}. Take \(d=c\) and
\(w(k)=1\) in Lemma~\ref{lem:slope}, and choose any prime
\(p\mid c\). The hypothesis \(\gcd(c,ab)=1\) gives \(p\nmid ab\).
Outside a set of density zero, the \(p\)-part of the denominator
of the truncation is at least
\[
 p^{\frac12\sigma_p(c)\log n}
   =n^{\frac12\sigma_p(c)\log p}.
\]
Proposition~\ref{prop:transfer} applies with one group and \(q_0=0\),
proving normality.

The mechanism can be read directly from the series. Between insertions,
the shifted truncations follow rational radix orbits. Their denominators
become large at almost every position, and the orbit estimate turns
that growth into the correct frequency of every digit word.

\section{Why rational combinations remain normal}

We turn to Theorem~\ref{thm:arithmetic}. Discard zero coefficients.
The key is to group parameters that are powers of the same integer
\emph{before} comparing their denominators.

\subsection*{First collect the exact overlaps}

Every integer \(c\ge2\) has a unique expression
\(c=d^r\), where \(r\ge1\) and \(d\) is not a proper integer
power. In its prime factorization, \(r\) is the greatest common
divisor of the exponents. Two integers have equal positive powers
exactly when their corresponding \(d\)'s agree.

Group together all parameters with the same \(d\). Their contribution
to \eqref{eq:main-combination} is
\begin{equation}\label{eq:grouping}
 Y_d=\sum_{i,k\ge0}\frac{w_d(k)}{a^i d^k b^{a^i d^k}},
 \qquad
 w_d(k)=\sum_{r\in E_d}\lambda_{d,r}\,\mathbf1_{r\mid k}.
\end{equation}
Here \(E_d\) is the finite set of powers present in the group, and
\(r\mid0\) as usual. The sequence \(w_d\) is periodic, with period
dividing \(\lcm(E_d)\).

It is also nonzero. Choose the smallest \(r\in E_d\) with nonzero
coefficient and evaluate at \(k=r\). Any other contributing exponent
must be a proper divisor of \(r\), whose coefficient is zero by the
choice of \(r\). Hence \(w_d(r)=\lambda_{d,r}\ne0\).

For instance, \(\xi_{b,d;a}-\xi_{b,d^2;a}\) has weight one at odd
\(k\) and zero at even \(k\). Grouping makes this exact cancellation
visible, and Lemma~\ref{lem:slope} applies to the surviving periodic
weight.

\subsection*{Then one prime separates the groups}

Let \(S_{d,n}\) be the truncation of \(Y_d\) at insertion position
\(n\). Choose a prime \(p\) dividing at least one active \(d\).
It divides neither \(a\) nor \(b\). Associate to each group the slope
\[
 \sigma_d=\frac{v_p(d)}{\log d},
\]
which is zero if \(p\nmid d\). Distinct groups cannot have equal
positive slopes: equality would give
\[
 d_1^{v_p(d_2)}=d_2^{v_p(d_1)},
\]
forcing the same primitive integer in both groups.

There is therefore a unique group \(d_*\) of largest positive slope.
Choose \(0<\theta<\sigma_{d_*}\) with
\(\sigma_d<\theta\) for every other group.
By Lemma~\ref{lem:slope}, outside a set of density zero,
\[
 v_p(S_{d_*,n})<-\theta\log n.
\]
Every other group satisfies the unconditional bound
\[
 v_p(S_{d,n})\ge-\sigma_d\log n-O(1).
\]
When \(p\nmid d\), this simply says that its \(p\)-denominator is
bounded: only the finitely many rational weights can contribute such a
factor. The rational constant \(q_0\) has a fixed valuation as well.
For all sufficiently large nonexceptional \(n\), the \(d_*\) group
has strictly smaller valuation than every other contribution.
Rule~\eqref{eq:valuation-rule} gives
\[
 v_p\!\left(q_0+\sum_d S_{d,n}\right)
       =v_p(S_{d_*,n})<-\theta\log n.
\]
Thus the \(p\)-part of the denominator of the complete truncation
exceeds \(n^{\theta\log p}\) at a density-one set of positions.
Proposition~\ref{prop:transfer} proves that
\(q_0+\sum_dY_d\) is normal in base \(b\), completing the proof
of Theorem~\ref{thm:arithmetic}.\hfill\(\square\)

\Needspace{8\baselineskip}
\begin{corollary}\label{cor:independence}
For fixed \(a,b\ge2\), the number one and all the values
\(\xi_{b,c;a}\), with \(c\ge2\) and \(\gcd(c,ab)=1\), are
linearly independent over \(\Q\). The Xi values span an
infinite-dimensional rational vector space whose every nonzero element
is normal in base \(b\).
\end{corollary}

\begin{proof}
A nontrivial rational relation involving a Xi value would make the
normal number in Theorem~\ref{thm:arithmetic} equal to zero. This is
impossible: rational numbers have eventually periodic expansions and
are not normal. There are infinitely many eligible parameters \(c\).
\end{proof}

The arithmetic theorem fixes both \(a\) and \(b\). Within that
family it allows every finite rational combination, including negative
coefficients and overlapping parameter powers. Its proof identifies
what makes this possible: periodic cancellations preserve a definite
denominator growth rate, and distinct groups have different rates at
a suitable prime.

\section*{The accompanying Lean proof}

Theorems~\ref{thm:normal} and~\ref{thm:arithmetic} are formalized
with their stated hypotheses. In the accompanying source package,
\texttt{XiFamilyChallenge.lean} defines the distinct-index series and
normality through overlapping word frequencies. It imports only
mathlib, independently of the proof development.
\texttt{XiFamilySolution.lean} proves both statements, and
\texttt{XiFamilyStatement.lean} establishes correspondence with the
definitions used internally.

The formal proof follows the same denominator argument. Its analytic
part proves sufficient rational-orbit estimates by repeated differencing,
so it does not assume Vandehey's theorem as an axiom. The complete
source build and the endpoint axiom checks passed with Lean
4.34.0-rc2 and pinned mathlib dependencies. There are no unfinished
proofs or added axioms. Convergence, the double-series identity, and
the Stoneham specialization are included in the development.

\begin{thebibliography}{9}
\small
\bibitem{Stoneham}
R.~G. Stoneham, On absolute \((j,\varepsilon)\)-normality in the rational
fractions with applications to normal numbers,
\emph{Acta Arithmetica} \textbf{22} (1973), 277--286.
\href{https://doi.org/10.4064/aa-22-3-277-286}{doi:10.4064/aa-22-3-277-286}.

\bibitem{BaileyCrandall}
D.~H. Bailey and R.~E. Crandall, Random generators and normal numbers,
\emph{Experimental Mathematics} \textbf{11} (2002), 527--546.
\href{https://www.davidhbailey.com/dhbpapers/bcnormal-em.pdf}{Author's copy}.

\bibitem{Dibag}
I. Dibag, Generalization of the von Staudt--Clausen theorem,
\emph{Journal of Algebra} \textbf{125} (1989), 519--523.
\href{https://doi.org/10.1016/0021-8693(89)90180-4}{doi:10.1016/0021-8693(89)90180-4}.

\bibitem{Kano}
H. Kano, General constructions of normal numbers of Korobov type,
\emph{Osaka Journal of Mathematics} \textbf{30} (1993), 909--919.
\href{https://doi.org/10.18910/10356}{doi:10.18910/10356}.

\bibitem{KanoShiokawa}
H. Kano and I. Shiokawa, Rings of normal and nonnormal numbers,
\emph{Israel Journal of Mathematics} \textbf{84} (1993), 403--416.
\href{https://www.math.keio.ac.jp/wp-content/uploads/2022/03/91006.pdf}{Institutional preprint}.

\bibitem{Vandehey}
J. Vandehey, Differencing methods for Korobov-type exponential sums,
\emph{Journal d'Analyse Math\'ematique} \textbf{138} (2019), 405--439.
\href{https://doi.org/10.1007/s11854-019-0038-2}{doi:10.1007/s11854-019-0038-2}.
\href{https://arxiv.org/abs/1606.07911}{arXiv:1606.07911}.

\bibitem{Weyl}
H. Weyl, \"Uber die Gleichverteilung von Zahlen mod.\ Eins,
\emph{Mathematische Annalen} \textbf{77} (1916), 313--352.
\href{https://doi.org/10.1007/BF01475864}{doi:10.1007/BF01475864}.
\end{thebibliography}
\end{document}
